\documentclass[11pt]{article}
\usepackage[a4paper,margin=2.7cm]{geometry}
\usepackage{amsmath,amssymb,amsthm,mathtools}
\usepackage{enumitem}
\usepackage{xcolor}
\usepackage{microtype}
\usepackage[colorlinks,linkcolor=blue!50!black,citecolor=blue!50!black,urlcolor=blue!50!black]{hyperref}
\hypersetup{pdftitle={Failure of a local continuation principle in Enflo's invariant subspace argument},
pdfsubject={An explicit model obstruction with an abstract Lean proof}}

\newtheorem{theorem}{Theorem}
\newtheorem{lemma}[theorem]{Lemma}
\newtheorem{proposition}[theorem]{Proposition}
\theoremstyle{definition}
\newtheorem{definition}[theorem]{Definition}
\newtheorem{remark}[theorem]{Remark}

\newcommand{\C}{\mathbb{C}}
\newcommand{\R}{\mathbb{R}}
\newcommand{\N}{\mathbb{N}}
\newcommand{\ip}[2]{\langle #1,#2\rangle}
\newcommand{\norm}[1]{\lVert #1\rVert}
\newcommand{\abs}[1]{\lvert #1\rvert}
\newcommand{\re}{\operatorname{Re}}
\newcommand{\cC}{\mathcal{C}}
\newcommand{\ha}{\hat a}
\newcommand{\hb}{\hat b}
\newcommand{\one}{\mathbf 1}
\newcommand{\rhoN}{\varrho}

\title{Failure of a local continuation principle in Enflo's
invariant subspace argument}
\author{}
\date{Draft, November 2026}

\begin{document}
\maketitle

\begin{abstract}
In arXiv:2305.15442v2, P.~Enflo proposes a proof that every bounded operator on a Hilbert space
has a non-trivial closed invariant subspace. The argument constructs approximate non-cyclic
vectors and claims that, for operators of ``type 1'', a constrained iterative construction
drives the approximation error to zero while the vectors converge in norm. We show that the
justification of this convergence step cannot be correct. For multiplication by the variable on
$L^2[0,\kappa]$, which is of type 1 and satisfies the other structural hypotheses used, we
exhibit explicit states satisfying the admissibility conditions extracted below, including
the non-degeneracy conditions, with arbitrarily small defect, from which no sequence obeying
the confinement and distance-drift constraints can have defect tending to zero. A supplementary
construction connects the starting vector to such states while preserving these state
conditions. It does not establish the degree or coefficient bounds needed to validate the
paper's first-order step formulas.
The obstruction is an elementary inequality, valid for every bounded operator, which we have
checked in Lean. This does not refute the theorem, and it says nothing about operators without
invariant subspaces; it shows that the published argument does not establish its key step.
\end{abstract}

\section{Introduction}

Let $H$ be a complex Hilbert space. Whether every bounded operator on a separable
infinite-dimensional $H$ has a closed invariant subspace other than $\{0\}$ and $H$ is the
invariant subspace problem. This note examines the argument in version 2 of Enflo's
preprint \cite{E}, whose source is included in the accompanying repository. Our claim concerns
the continuation inference in that version, independently of the status of other approaches
to the invariant subspace problem.

The argument of \cite{E} has three parts.
\begin{enumerate}[leftmargin=2em]
\item A reformulation. For a unit vector $x_0$, one seeks $v\notin\{0,x_0\}$ whose orbit
$\{T^jv\}$ is orthogonal to $x_0-v$. Extremal power series produce vectors $v$ for which all
the numbers $\ip{T^jv}{x_0-v}$ are bounded by a single non-negative number, which we call the
\emph{defect}. A norm limit of such vectors with defect tending to $0$ is non-cyclic.
\item A ``Main Construction'' (MC) that modifies the current state step by step so as to
decrease the defect, using first-order perturbation formulas.
\item A convergence step. Operators are divided into two types. For type 1 it is asserted that
one can add a linear constraint to the construction which confines the vectors to a small ball
and still continue to defect $0$; repeating this with tighter constraints gives norm
convergence. One branch of the treatment of type 2 operators falls back on the same step.
\end{enumerate}
Part 1 is correct, and the device in part 3 for confining the vectors is correct and
attractive (Lemma~\ref{lem:confine}). Our subject is the assertion in part 3 that the
construction can be continued under the constraint. In \cite{E} this is derived from a
quantitative linear-independence claim, justified by the type 1 property.

\subsection*{Results}
We work with $T=$ multiplication by $\lambda$ on $L^2[0,\kappa]$. This operator has many
invariant subspaces, so it is not a counterexample to anything. It is a test of the reasoning:
it is of type 1, and every consequence of the non-existence of invariant subspaces that the
type 1 argument of \cite{E} invokes holds for it (Section~\ref{sec:discussion} lists them). If
the reasoning were valid, its conclusion would hold in this operator at every state satisfying
the hypotheses.
\begin{itemize}[leftmargin=2em]
\item Theorem~\ref{thm:main}: there are states satisfying
Definition~\ref{def:admissible}, with arbitrarily small defect, at which Claim (C) below
fails. Explicit extremal series and a compactness argument for the second coefficient
sequence establish the required non-degeneracy conditions.
\item Theorem~\ref{thm:path}: these states are connected to the starting vector of \cite{E} by
a finite chain of steps $y\mapsto y+\sum_jr_jT^jy$ along which the defect decreases strictly
and the admissibility and distance-drift conditions hold. This concerns the states, not
control of the coefficients or degrees of the steps (Remark~\ref{rem:path}).
\end{itemize}

\subsection*{Scope}
We do not claim that the theorem of \cite{E} is false, nor that no argument of this general
shape can work. The obstruction uses a structure that an operator without invariant subspaces
cannot have; Section~\ref{sec:discussion} explains why this is unavoidable in any test of the
argument and what it leaves open. The steps of the chain in Theorem~\ref{thm:path} are small in
norm but are not small perturbations in the sense of coefficient size, and we make no claim
about what a particular rule for choosing MC steps would do. The starting residual level is
$O(t_0^2)$ above $1/2$, rather than in the interval ending at $1/2$ prescribed by
\cite{E}, Lemma~1; see the discussion after Lemma~\ref{lem:x0}.

The logical target is the implication from the listed local state conditions to continued
descent under confinement. Theorem~\ref{thm:main} disproves this implication in the model.
Interpreting it as a criticism of \cite{E} requires the hypothesis audit in
Section~\ref{sec:discussion}; we do not replace the global assumption of transitivity by
a claim that the model is transitive.

\subsection*{Conventions}
$\ip\cdot\cdot$ is linear in the first variable. $T$ is a bounded operator with
$\norm T=\kappa\le\frac12$ (in \cite{E}, $\kappa=10^{-20}$), injective, with dense range and not
surjective; \cite{E} reduces to this case. $\ell^2=\ell^2(\N_0)$, with shift $S$.

\section{The extremal framework}

For $y\in H$ let $V_y:\ell^2\to H$, $V_ya=\sum_{j\ge0}a_jT^jy$; this is bounded since
$\norm T<1$, and $TV_y=V_yS$. Fix $x_0$ with $\norm{x_0}=1$.

\begin{definition}
Let $y$ be cyclic and $0<\varepsilon<1$. The \emph{extremal series} of $y$ at level
$\varepsilon$ is the $a\in\ell^2$ of least norm with $\norm{x_0-V_ya}\le\varepsilon$, and
$v=V_ya$ is the \emph{extremal image}. For $v\in H$ put
\[
\tau(v)=\ip{v}{x_0-v},\qquad\Delta(v)=\sup_{j\ge0}\abs{\ip{T^jv}{x_0-v}} .
\]
\end{definition}

\begin{lemma}[\cite{E}, (1)--(17)]\label{lem:framework}
\begin{enumerate}[label=\textup{(\roman*)}]
\item The extremal series exists and is unique, $\norm{x_0-v}=\varepsilon$, and there is $t>0$
with $ta=V_y^*(x_0-v)$. Consequently $\ip{T^jv}{x_0-v}=t\ip{S^ja}{a}$ for all $j$, so
$\tau(v)=t\norm a^2>0$ and $\Delta(v)=\tau(v)$. Moreover $\norm v\le1$ and
$x_0-v=t(V_yV_y^*+t)^{-1}x_0$.
\item Conversely, if $y$ is cyclic, $0\ne a\in\ell^2$, $t>0$ and $ta=V_y^*(x_0-V_ya)$, then
$a$ is the extremal series of $y$ at level $\norm{x_0-V_ya}$.
\end{enumerate}
\end{lemma}
\begin{proof}
This is Tikhonov regularisation. For $t>0$ the functional
$F_t(b)=\norm{x_0-V_yb}^2+t\norm b^2$ is strictly convex with a unique minimiser, characterised
by the Euler--Lagrange equation $tb=V_y^*(x_0-V_yb)$. If $a$ satisfies it, and
$\norm{x_0-V_yb}\le\norm{x_0-V_ya}$, $\norm b\le\norm a$, then $F_t(b)\le F_t(a)$ and $b=a$;
this is (ii). The residual norm is continuous and strictly increasing in $t$, from $0$ to $1$,
which gives existence in (i). Then
$\ip{T^jv}{x_0-v}=\ip{V_yS^ja}{x_0-v}=\ip{S^ja}{V_y^*(x_0-v)}=t\ip{S^ja}{a}$, and
$1=\norm v^2+\varepsilon^2+2\tau(v)$. Finally
$(V_yV_y^*+t)(x_0-v)=tV_ya+t(x_0-v)=tx_0$.
\end{proof}

We call $\tau(v)$ the \emph{defect} of an extremal image; it is $\varepsilon\theta$ in \cite{E}.
If $\Delta(v)=0$ and $v\notin\{0,x_0\}$, the closed span of the orbit of $v$ is a non-trivial
invariant subspace. Since $\Delta$ is norm-continuous, it suffices to find extremal images
$v_n$ that converge in norm, with $\tau(v_n)\to0$ and $\norm{x_0-v_n}$ bounded away from $0$
and $1$. Norm convergence is the whole difficulty: $\Delta$ is not weakly continuous.

\subsection*{Balanced coefficients}
\cite{E} normalises $y$ so that $t=1$: with $\hat y=t^{-1/2}y$ and
$R=(I+V_{\hat y}V_{\hat y}^*)^{-1}$ (written $[\;]^{-1}$ there) one has $Rx_0=x_0-v$ by
Lemma~\ref{lem:framework}(i). It uses the two sequences
$a_m=\ip{Rx_0}{T^m\hat y}$ and $b_m=\ip{R^2x_0}{T^m\hat y}$. In terms of an arbitrary
normalisation $(y,a,t)$ these are
\begin{equation}\label{eq:balanced}
\ha=t^{1/2}a,\qquad \hb=t^{3/2}\,(t+V_y^*V_y)^{-1}a ,
\end{equation}
because $\hb=V_{\hat y}^*R(x_0-v)=(I+V_{\hat y}^*V_{\hat y})^{-1}V_{\hat y}^*(x_0-v)$ and
$V_{\hat y}^*(x_0-v)=t^{1/2}a$. Note $\norm{\ha}^2=\tau(v)$.

\subsection*{The confinement device}

\begin{lemma}\label{lem:confine}
Let $v,v'$ be extremal images, $y'$ a unit vector with
$\norm{v'/\norm{v'}-y'}\le\vartheta$ and $w\in H$ with $\norm{w-y'}\le\vartheta$. If
$\ip{x_0-v}{w}=\ip{x_0-v'}{w}$, then
\[
\norm{v-v'}^2\le2\tau(v')+8\vartheta+\bigl|\norm{x_0-v'}^2-\norm{x_0-v}^2\bigr| .
\]
\end{lemma}
\begin{proof}
With $r=x_0-v$, $r'=x_0-v'$ one has the identity
$\norm{v-v'}^2=\norm{r-r'}^2=2\re\ip{v'}{r}-2\tau(v)+\norm{r'}^2-\norm r^2$, and
$\tau(v)\ge0$. As $\norm{v'},\norm r,\norm{r'}\le1$ and $\norm{v'/\norm{v'}-w}\le2\vartheta$,
$\re\ip{v'}{r}\le\norm{v'}\re\ip wr+2\vartheta=\norm{v'}\re\ip w{r'}+2\vartheta
\le\re\ip{v'}{r'}+4\vartheta=\tau(v')+4\vartheta$.
\end{proof}

So holding a single inner product fixed confines later states to a ball around $v'$ whose
radius is small when the defect of $v'$ is small. This is how \cite{E} proposes to obtain
norm convergence.

\subsection*{The hypotheses on a state}
\cite{E} calls $T$ of \emph{type 1} if there is a unit vector $u_0$ such that for every
$n\ge1$ there is $\delta_n>0$ with: every $y\ne0$ with
$\re\ip{y}{u_0}\ge\frac1{100}\norm y$ has $\abs{\ip{T^jy}{y}}\ge\delta_n\norm y^2$ for some
$j\ge n$. For such $T$ it takes $x_0=\frac{\sqrt3}2u_0+\frac12u_1$ with $u_1\perp u_0$ a unit
vector and $\norm{T^*u_1}$ smaller than a number $(\varepsilon\theta)_0$ fixed in advance, and
$y_0=\frac{\sqrt3}2u_0$ as the starting vector, so that $\norm{x_0-y_0}=\frac12$. The following
collects what \cite{E} arranges before, and maintains during, the MC.

\begin{definition}\label{def:admissible}
Let $T$ be of type 1 with vector $u_0$. An extremal image $v=V_ya$, with parameter $t$ and
defect $\tau=\tau(v)$, is \emph{admissible} if
\begin{enumerate}[label=\textup{(A\arabic*)},leftmargin=3em]
\item $y$ is cyclic, $\re\ip y{u_0}\ge\frac1{100}\norm y$, and $0.3<\norm{x_0-v}<0.7$
\textup{(\cite{E}, p.~4 and (19))};
\item the leading coefficient dominates: $\abs{a_0}^2\ge100\sum_{j\ge1}\abs{a_j}^2$
\textup{(\cite{E}, (22))};
\item the tail is not too small: $\sum_{j\ge1}\abs{a_j}^2\ge\tau^{15}\norm a^2$
\textup{(\cite{E}, (25) and (33))};
\item there are four distinct indices $j_1,j_2,j_3,j_4\ge2$ with
$\abs{\ha_{j_1}},\abs{\ha_{j_2}},\abs{\hb_{j_3}},\abs{\hb_{j_4}}>\tau^{30}$
\textup{(\cite{E}, Remark after (27))}.
\end{enumerate}
\end{definition}

Some remarks on how this matches the text of \cite{E}. (1) The cone condition there is
$\ip{y/\norm y}{u_0}\ge\frac1{100}$; with the real part the cone is larger, so type 1 in
our sense implies type 1 in the sense of \cite{E}, and the type 1 inequalities for $y$ hold
by (A1). (2) \cite{E} states the range in (A1) once for $\varepsilon=\norm{x_0-v}$ and once
(after (37)) for $\ip{Rx_0}{x_0}=\varepsilon^2+\tau$. The second cannot be meant, since it
fails at the paper's own start, where $\varepsilon=\frac12$. (3) (A3) is the stronger of the
two forms in \cite{E}: (25) is stated in the normalisation $a_0\approx1$, and (33) has a
factor $\frac12$. (4) In (A4) we read $a_m,b_m$ in the balanced normalisation in which
\cite{E} defines them; the states below satisfy (A4) in the normalisation $\norm y\approx1$
as well. (5) \cite{E}, (33) also contains an upper bound for the tail, of the form
$\gamma'(\tau)\norm a^2$ with an unspecified function $\gamma'\to0$ coming from its
Lemma~2. It is described there as automatic and we do not make it a hypothesis; for all
states below the tail is at most $\tau^2\norm a^2$.

\subsection*{The claim under examination}
At a stage $n_0'$ of the MC where the defect is small, with $y'=y_{n_0'}/\norm{y_{n_0'}}$, the
paper sets
\[
w_{00}=y'+\delta_2^{100}\,T^{*j_1}y'+\delta_2^{100}\,T^{*j_2}y'
\]
for certain indices $j_1,j_2$ supplied by the type 1 property, and requires all later steps to
keep $\ip{x_0-v_n}{w_{00}}$ fixed. It then asserts (\cite{E}, text after (46)) that the vectors
governing the constrained steps are
\begin{quote}
``$\delta_2$-linearly independent to arbitrarily small $\varepsilon\theta$'s, since with
the condition
\[\langle ch\,[\;]^{-1}_nx_0,w_{00}\rangle=0\]
the room for variation of $[\;]^{-1}_nx_0$ is less than $\delta_2^{40}$'',
\end{quote}
which ``will allow MC to be used to arbitrarily small $\varepsilon\theta$'s''. Earlier,
\cite{E} (45) bounds the change of $\ip{Rx_0}{x_0}=\norm{x_0-v}^2+\tau(v)$ during the
construction by $10$ times the initial defect. In our notation:

\begin{quote}
\textbf{Claim (C).} \emph{Let $T$ be of type 1 and let $v'$ be an admissible extremal image
with sufficiently small defect. Then there are extremal images $v_n$ with $\tau(v_n)\to0$,
$\ip{x_0-v_n}{w_{00}}=\ip{x_0-v'}{w_{00}}$ and
$\bigl|\norm{x_0-v_n}^2-\norm{x_0-v'}^2\bigr|\le11\,\tau(v')$.}
\end{quote}

In \cite{E} the stage $n_0'$ is ``suitable'', and ``we can wait until $(\varepsilon\theta)_n$
is arbitrarily small''; no criterion of suitability other than small defect is given, and the
same step is repeated at arbitrarily late stages with $w_{01},w_{02},\dots$. We therefore
state the claim for every admissible state of small defect. Section~\ref{sec:reach} addresses
the possibility that the states the MC visits are special.

The only reasons given are the type 1 inequalities, used to argue linear independence. Since
$\norm{w_{00}-y'}\le2\delta_2^{100}$, Lemma~\ref{lem:confine} shows that the $v_n$ of Claim (C)
all lie in a small ball around $v'$. So Claim (C) can hold only if that ball contains vectors
of arbitrarily small defect.

\section{A no-continuation inequality}

\begin{theorem}\label{thm:nogo}
Let $T$ be a bounded operator on $H$, $W$ an isometry of $H$, $x_0\in H$,
$v=\frac12(x_0+Wx_0)$ and $C=\frac12(W+W^*)$.
\begin{enumerate}[label=\textup{(\roman*)}]
\item For every $u\in H$,
\[
\norm{u-v}^2\ \ge\ \tfrac14\bigl(\norm{x_0}^2-\norm{Cx_0}^2\bigr)
+\tfrac12\re\ip{(W-W^*)u}{x_0-u}.
\]
\item Suppose $W=\sum_kc_kT^k$ and $W^*=\sum_kc_k'T^k$, convergent in operator norm, with
$\sum_k\abs{c_k}\le K$ and $\sum_k\abs{c_k'}\le K$. Then for every $u\in H$,
\[
\Delta(u)\ \ge\ \frac1K\Bigl(\tfrac14\bigl(\norm{x_0}^2-\norm{Cx_0}^2\bigr)-\norm{u-v}^2\Bigr).
\]
\item If $W$ is unitary and commutes with $T$, then
$\ip{T^jv}{x_0-v}=\frac14\ip{T^j(W-W^*)x_0}{x_0}$ for all $j\ge0$.
\end{enumerate}
\end{theorem}
\begin{proof}
(i) Let $B=-\frac14(W-W^*)$, so $B^*=-B$. Then $\re\ip{Bu}{u}=0$ and
$\re\ip{Bu}{x_0}=-\re\ip{u}{Bx_0}$, whence for every $u$
\[
\norm{u-v}^2+2\re\ip{Bu}{x_0-u}=\norm{u-(v+Bx_0)}^2+\norm v^2-\norm{v+Bx_0}^2
\ \ge\ \norm v^2-\norm{v+Bx_0}^2 .
\]
Now $v+Bx_0=\frac12x_0+\frac12Cx_0$. Using $\norm{Wx_0}=\norm{x_0}$ and
$\re\ip{Cx_0}{x_0}=\re\ip{Wx_0}{x_0}$,
$\norm v^2=\frac12\norm{x_0}^2+\frac12\re\ip{Wx_0}{x_0}$ and
$\norm{v+Bx_0}^2=\frac14\norm{x_0}^2+\frac12\re\ip{Wx_0}{x_0}+\frac14\norm{Cx_0}^2$.
(ii) $\abs{\ip{Wu}{x_0-u}}\le\sum_k\abs{c_k}\abs{\ip{T^ku}{x_0-u}}\le K\Delta(u)$ and likewise
for $W^*$; insert in (i). (iii) Expand, using $\ip{Wa}{Wb}=\ip ab$ and $WT^j=T^jW$.
\end{proof}

Part (ii) says: if $x_0$ is not nearly fixed by $C$, then every vector in a ball of definite
radius around $v$ has defect bounded below. Part (iii) says that $v$ itself has small defect
when $W-W^*$ is small against the orbit of $x_0$. Both can hold at once.

Theorem~\ref{thm:nogo} has been formalised in Lean~4 with Mathlib (file
\texttt{Oscillating.lean} of the accompanying project); parts (i), (ii), (iii) are the
theorems
\begin{center}
\texttt{oscillating\_lower\_bound},\quad\texttt{oscillating\_ball\_bound},\quad
\texttt{oscillating\_defect}.
\end{center}

\section{The model operator and its extremal images}\label{sec:model}

Let $0<\kappa\le\frac12$, $H=L^2[0,\kappa]$ with Lebesgue measure and
$(Tf)(\lambda)=\lambda f(\lambda)$. Let $u_0=\kappa^{-1/2}\one$ and
$k=\frac{\sqrt3}2\kappa^{-1/2}$, so that $y_0=\frac{\sqrt3}2u_0$ is the constant function $k$,
and $k^2\kappa=\frac34$.

\begin{lemma}\label{lem:model}
\begin{enumerate}[label=\textup{(\roman*)}]
\item $T$ is injective with dense range and not surjective; $f$ is cyclic iff $f\ne0$ a.e.
\item $T$ is of type 1 with the vector $u_0$ (indeed with any unit vector).
\item If $d\ne0$ and $\abs{c_j}\le1$ for $j\ge1$, then $d-\sum_{j\ge1}c_jT^j$ is injective.
\end{enumerate}
\end{lemma}
\begin{proof}
(i) is standard (Weierstrass). (ii) Pick $\eta>0$ with
$\norm{u_0\one_{[0,\eta]}}\le\frac1{200}$. If $\norm y=1$ and
$\re\ip y{u_0}\ge\frac1{100}$ then $\norm{y\one_{(\eta,\kappa]}}\ge\frac1{200}$, so
$\ip{T^jy}{y}=\int\lambda^j\abs y^2\ge\eta^j/40000$ for all $j$. (iii) The function
$d-\sum_jc_j\lambda^j$ is real-analytic and not identically zero, hence non-zero a.e.
\end{proof}

\subsection{\texorpdfstring{The choice of $x_0$}{The choice of x0}}
Let $L(z)=-z^{-1}\log(1-z)=\sum_{j\ge0}z^j/(j+1)$, so $1\le L(z)\le1+z$ for $0\le z\le\frac14$.
In this section ``$t_0$ small'' means smaller than a constant depending only on $\kappa$.

\begin{lemma}\label{lem:x0}
For every small $t_0>0$ there are $\mu\in(\frac12,1)$ and $\sigma\in(0,8t_0^2/k^2)$ such that,
with $h=t_0\mu/(k\sigma)$,
\[
x_0\;=\;k\mu\,L(\sigma\lambda)+h\,\one_{[0,\sigma]}
\]
is a unit vector with $\ip{x_0}{u_0}=\frac{\sqrt3}2$. Put $u_1=2(x_0-y_0)$ and
$\varepsilon_0=h\sigma^{1/2}$. Then $u_1$ is a unit vector orthogonal to $u_0$,
$x_0=\frac{\sqrt3}2u_0+\frac12u_1$, $\norm{T^*u_1}\le\frac23t_0+O(t_0^2)$,
$\abs{\mu^{-1}-1-\frac43t_0}=O(t_0^2)$ and $\varepsilon_0^2=\frac14+O(t_0^2)$. Moreover the
extremal image of $y_0$ with parameter $t_0$ is
\[
v_0=k\mu L(\sigma\lambda)=x_0-h\one_{[0,\sigma]},\qquad\text{with series }
a^{(0)}_j=\frac{\mu\,\sigma^j}{j+1},
\]
so $\norm{x_0-v_0}=\varepsilon_0$ and
$\tau(v_0)=t_0\mu^2\Lambda(\sigma)$, where $\Lambda(\beta)=\sum_{j\ge0}\beta^{2j}/(j+1)^2$.
\end{lemma}
\begin{proof}
The condition $\ip{x_0}{u_0}=\frac{\sqrt3}2$ reads $k\mu I(\sigma)+h\sigma=k\kappa$ with
$I(\sigma)=\int_0^\kappa L(\sigma\lambda)\,d\lambda\in[\kappa,\kappa(1+\sigma\kappa)]$, and
$h\sigma=t_0\mu/k$, so it determines $\mu=\mu(\sigma)=\kappa/(I(\sigma)+t_0k^{-2})$,
continuous in $\sigma$. With this $\mu$,
\[
\norm{x_0}^2=\mu^2\Bigl(k^2\!\int_0^\kappa\!L(\sigma\lambda)^2d\lambda
+\frac{2t_0}\sigma\!\int_0^\sigma\!L(\sigma\lambda)\,d\lambda+\frac{t_0^2}{k^2\sigma}\Bigr)
\]
is continuous in $\sigma$, tends to $\infty$ as $\sigma\to0$, and at $\sigma=8t_0^2/k^2$ is at
most $(1+\frac43t_0)^{-2}(\frac34+\frac18+O(t_0))<1$. So there is a root
$\sigma<8t_0^2/k^2$. The asymptotics follow by expanding the two equations. As
$\norm{x_0}=1$ and $\ip{x_0}{y_0}=\frac34=\norm{y_0}^2$, $u_1$ is a unit vector orthogonal to
$u_0$; and $u_1=2k(\mu L(\sigma\lambda)-1)+2h\one_{[0,\sigma]}$ gives
$\norm{Tu_1}\le2k\,\bigl(1-\mu+O(t_0^2)\bigr)\kappa^{3/2}/\sqrt3+2h\sigma^{3/2}
\le\frac43\kappa t_0+O(t_0^2)$; and $T^*=T$.

For the last statement, $V_{y_0}a^{(0)}=k\sum_ja^{(0)}_j\lambda^j=v_0$ and
\[
\ip{x_0-v_0}{T^jy_0}=kh\!\int_0^\sigma\!\lambda^j\,d\lambda=\frac{t_0\mu\,\sigma^{j}}{j+1}
=t_0a^{(0)}_j ,
\]
so Lemma~\ref{lem:framework}(ii) applies, and $\tau(v_0)=t_0\norm{a^{(0)}}^2$.
\end{proof}

From now on $t_0$ is small and $x_0,\mu,\sigma,h,\varepsilon_0$ are as in Lemma~\ref{lem:x0}.
Thus $x_0$ is of the form prescribed in \cite{E}, and $v_0$ is the extremal image of the
paper's starting vector $y_0$ at the level $\varepsilon_0=\frac12+O(t_0^2)$; its defect
$\tau(v_0)=t_0(1+O(t_0))$ is as small as we please and exceeds $\norm{T^*u_1}$, so the
requirement $\norm{T^*u_1}<(\varepsilon\theta)_0$ of \cite{E} holds for any
$(\varepsilon\theta)_0\ge\tau(v_0)$. Two differences from \cite{E} should be noted.
\cite{E}, Lemma~1 takes a level in $(\frac12-10^{-5}(\varepsilon\theta)_0,\frac12]$; ours is
$O(t_0^2)$ above $\frac12$. And \cite{E} does not start the MC at the extremal image of $y_0$
itself but after a preparation (its Lemma~2 and Cases I, II) whose purpose is to obtain
(A2)--(A4); here $v_0$ already has these properties (Theorem~\ref{thm:path}), so no
preparation is needed. We record that $h\ge k$, that $h\le x_0\le(1+9t_0)h$ on $[0,\sigma]$, and that
$\frac12k\le x_0\le1.01k$ on $(\sigma,\kappa]$.

\subsection{A formula for extremal images}
In the model every extremal image can be written down from a single function.

\begin{lemma}\label{lem:explicit}
Let $\psi\in L^\infty[0,\kappa]$ be complex-valued with $t:=\int x_0^2\re\psi\,d\lambda>0$.
Put
\[
a_j=\frac1t\int_0^\kappa\lambda^jx_0^2\,\overline\psi\,d\lambda,\qquad
A(\lambda)=\sum_{j\ge0}a_j\lambda^j,\qquad c=\psi A .
\]
Suppose that $A$ has no zero on $[0,\kappa]$ and that $\re c+(\operatorname{Im}c)^2\le\frac14$
a.e. Let $\varsigma:[0,\kappa]\to\{1,-1\}$ be measurable and
\[
g=\tfrac12+\varsigma\sqrt{\tfrac14-\re c-(\operatorname{Im}c)^2}+i\operatorname{Im}c,
\qquad v=x_0g,\qquad y=v/A .
\]
If $g\ne0$ a.e., then $y$ is cyclic, $v$ is the extremal image of $y$ at level
$\norm{x_0-v}$, with extremal series $a$ and parameter $t$; and $\tau(v)=t\norm a^2$.
\end{lemma}
\begin{proof}
The series for $A$ converges uniformly on $[0,\kappa]$ since
$\abs{a_j}\le\kappa^jt^{-1}\norm{\psi}_\infty$. One checks $g-\abs g^2=c$. Then
$V_ya=Ay=v$ and, as $x_0$ is real,
\[
\ip{x_0-v}{T^jy}=\int\lambda^jx_0^2\,\frac{(1-g)\overline g}{\overline A}\,d\lambda
=\int\lambda^jx_0^2\,\frac{\overline c}{\overline A}\,d\lambda
=\int\lambda^jx_0^2\,\overline\psi\,d\lambda=t\,a_j .
\]
Now apply Lemma~\ref{lem:framework}.
\end{proof}



\subsection{A family of states}\label{ssec:states}
For $0<\beta\le\sigma$, $0\le\alpha\le\alpha_{\max}:=1.54$ and $N\in\N$ let
$S(\beta,\alpha,N)$ be the data produced by Lemma~\ref{lem:explicit} from
\[
\psi=\begin{cases}\dfrac{\mu^2t_0}{\sigma}\,x_0^{-2}&\text{on }[0,\beta],\\[2ex]
\dfrac i2\,\dfrac{\sin(\alpha\sin N\lambda)}{L(\beta\lambda)}&\text{on }(\beta,\kappa],
\end{cases}
\qquad \varsigma=\begin{cases}-1&\text{on }[0,\beta],\\+1&\text{on }(\beta,\kappa].\end{cases}
\]
Then $t=t(\beta):=\mu^2t_0\beta/\sigma$ and
\begin{equation}\label{eq:aj}
a_j=\frac{\beta^j}{j+1}-\frac{i}{t}\,p_j,\qquad
p_j=\frac12\int_\beta^\kappa\lambda^jx_0^2\,\frac{\sin(\alpha\sin N\lambda)}{L(\beta\lambda)}\,
d\lambda\in\R .
\end{equation}
The first term of $a_j$ is the $j$-th moment of the uniform probability measure on
$[0,\beta]$ and sums to $L(\beta\lambda)$; the second is small when $N$ is large, by
cancellation. On $(\beta,\kappa]$ the state is, up to this small term, $x_0$ multiplied by
$\frac12(1+e^{i\alpha\sin N\lambda})$; on $[0,\beta]$ it is close to $0$.

We use Fej\'er's lemma in the following form: if $\mathcal F\subset L^1[0,\kappa]$ and
$\mathcal H\subset C(\R/2\pi\mathbb Z)$ are norm-compact, then
$\int f(\lambda)h(N\lambda)\,d\lambda\to\bar h\int f$ as $N\to\infty$, uniformly for
$f\in\mathcal F$, $h\in\mathcal H$, where $\bar h$ is the mean of $h$. Put
\[
\omega_N=\sup\Bigl\{\sum_{j\ge0}\abs{p_j}:0<\beta\le\sigma,\ 0\le\alpha\le\alpha_{\max}\Bigr\},
\qquad\rhoN=\rhoN(\beta,N)=\omega_N/t(\beta).
\]
Since $t\mapsto\sin(\alpha\sin t)$ has mean $0$ and $\abs{p_j}\le\frac12\kappa^j$, Fej\'er's
lemma gives $\omega_N\to0$. Let
\[
m(\beta)=\int_0^\beta x_0^2,\qquad
g_\beta=\tfrac12-\sqrt{\tfrac14-\tfrac{\mu^2t_0}{\sigma}L(\beta\lambda)x_0^{-2}},\qquad
D(\beta)=\int_0^\beta x_0^2(1-g_\beta)^2 ,
\]
\[
\varepsilon_\infty^2(\beta,\alpha)=D(\beta)+\tfrac12\bigl(1-J_0(\alpha)\bigr)\bigl(1-m(\beta)\bigr),
\qquad v_N^\alpha=\tfrac12\bigl(x_0+e^{i\alpha\sin(NT)}x_0\bigr),
\]
where $J_0$ is the Bessel function; recall that the mean of $\sin^2(\frac12\alpha\sin t)$ is
$\frac12(1-J_0(\alpha))$.

\begin{proposition}\label{prop:states}
Let $0<\beta\le\sigma$, $0\le\alpha\le\alpha_{\max}$ and $\rhoN\le10^{-5}$. Then
$S(\beta,\alpha,N)=(y,a,t,v)$ is well defined: $v$ is the extremal image of the cyclic vector
$y$, with series $a$ and parameter $t=t(\beta)$. Moreover:
\begin{enumerate}[label=\textup{(\roman*)}]
\item $\sum_j\abs{a_j-\beta^j/(j+1)}\le\rhoN$, $\abs{a_j}\ge\beta^j/(j+1)$, and
$\tau(v)=t(\beta)\bigl(\Lambda(\beta)+\delta\bigr)$ with $0\le\delta\le\rhoN^2$;
\item $\frac14\beta^2\le\sum_{j\ge1}\abs{a_j}^2\le\beta^2+\rhoN^2$;
\item $\abs y\ge\frac15k$ a.e., $\norm y^2\le2$, $\re\ip y{u_0}\ge\frac1{100}\norm y$, and
$\bigl\|v/\norm v-y/\norm y\bigr\|\le2(\sigma+\rhoN)$;
\item $\bigl|\norm{x_0-v}^2-\varepsilon_\infty^2(\beta,\alpha)\bigr|\le100\rhoN+\phi_N$,
where $\phi_N\to0$ as $N\to\infty$ uniformly in $(\beta,\alpha)$;
\item $\norm{v-v_N^\alpha}^2\le1.11\,m(\beta)+(100\rhoN)^2$, and
$m(\beta)\le0.31\,\beta/\sigma$.
\end{enumerate}
For fixed $N$ the map $(\beta,\alpha)\mapsto(y,a,v)$ is norm-continuous where it is defined,
and $S(\sigma,0,N)=(\mu y_0,\mu^{-1}a^{(0)},\mu^2t_0,v_0)$ for every $N$: the state at
$(\sigma,0)$ is the extremal image of $y_0$ of Lemma~\ref{lem:x0}.
\end{proposition}
\begin{proof}
By \eqref{eq:aj}, $A=L(\beta\lambda)(1+\zeta)$ with $\abs\zeta\le\rhoN$ on $[0,\kappa]$, since
$L\ge1$. So $A$ has no zeros and $0.99\le\abs A\le1.01$.

\emph{On $(\beta,\kappa]$.} Here $c=\frac i2s(1+\zeta)$ with $s=\sin(\alpha\sin N\lambda)$,
$\abs s\le\sin\alpha_{\max}$, and
\[
\tfrac14-\re c-(\operatorname{Im}c)^2=\tfrac14(1-s^2)+\epsilon,\qquad\abs\epsilon\le1.25\rhoN .
\]
As $\frac14\cos^2\alpha_{\max}>2.3\cdot10^{-4}$, the radicand is positive, and comparing with
$g_N:=\frac12(1+e^{i\alpha\sin N\lambda})=\frac12+\frac12\sqrt{1-s^2}+\frac i2s$ we get
$\abs{g-g_N}\le1.25\rhoN/(\frac12\cos\alpha_{\max})+\frac12\rhoN\le82\rhoN$. In particular
$\re g\ge\frac12$, $\abs g\le1.01$, $\re(g/A)\ge0.48$ and $\abs{g/A}\le1.04$. Also
$\abs{1-g}^2=1-\re g-\re c$ differs from
$\abs{1-g_N}^2=\sin^2(\frac12\alpha\sin N\lambda)$ by at most $82\rhoN$.

\emph{On $[0,\beta]$.} Here $c=c_\beta(1+\zeta)$ with
$c_\beta=\mu^2t_0\sigma^{-1}L(\beta\lambda)x_0^{-2}\in(0,5t_0]$, because $x_0\ge h$ and
$\mu^2t_0/(\sigma h^2)=\mu^2t_0/\varepsilon_0^2$. For $\abs c\le0.05$ the root
$g=\frac12-\sqrt{\frac14-\re c-(\operatorname{Im}c)^2}+i\operatorname{Im}c$ satisfies
$\abs{g-c}\le2.5\abs c^2$ and depends on $c$ with Lipschitz constant at most $3$. Hence the
radicand is positive, $g\ne0$, $\abs{g-g_\beta}\le\rhoN$, and
$y=x_0g/A=\mu^2t_0\sigma^{-1}x_0^{-1}(1+\zeta')$ with $\abs{\zeta'}\le0.11$; as
$\mu^2t_0/(\sigma h)=\mu k$ this gives $0.4k\le\abs y\le1.11k$ and $\re y\ge0$ there.

So Lemma~\ref{lem:explicit} applies. (i) follows from \eqref{eq:aj}, because $p_j$ is real:
$\abs{a_j}^2=\beta^{2j}/(j+1)^2+p_j^2/t^2$. (ii) follows from (i). For (iii), on
$(\beta,\kappa]$ we have $0.48x_0\le\abs y\le1.04x_0$ and $x_0\ge\frac12k$, so
$\abs y\ge\frac15k$ everywhere and $\norm y^2\le1.09+1.24k^2\sigma\le2$; and
$\ip y{u_0}=\kappa^{-1/2}\int y$ has real part at least
$0.48\,\kappa^{-1/2}\int_\sigma^\kappa x_0\ge0.2$. Finally $\norm{v-y}=\norm{(A-1)y}\le
(\sigma\kappa+\rhoN)\norm y$. For (iv) and (v) combine the two comparisons above:
$\abs{1-g}^2$ differs from $(1-g_\beta)^2$ on $[0,\beta]$ and from
$\sin^2(\frac12\alpha\sin N\lambda)$ on $(\beta,\kappa]$ by at most $82\rhoN$, and
$\phi_N:=\sup_{\beta,\alpha}\bigl|\int_\beta^\kappa x_0^2\sin^2(\frac12\alpha\sin N\lambda)
-\frac12(1-J_0(\alpha))(1-m(\beta))\bigr|\to0$ by Fej\'er's lemma; and
$\abs{g-g_N}\le1.05$, so $\abs{g-g_N}^2\le1.11$, on $[0,\beta]$, where $x_0^2\le1.2h^2$ and $h^2\sigma=\varepsilon_0^2$.

Continuity is clear from the formulas (the functions $\psi$ depend continuously on
$(\beta,\alpha)$ in $L^1(x_0^2d\lambda)$ and boundedly, so the $a_j$, and then $A$ uniformly,
and then $g$ in measure and boundedly). At $(\sigma,0)$ we have $p_j=0$,
$A=L(\sigma\lambda)$, $c=0$ and $g=1$ on $(\sigma,\kappa]$, while on $[0,\sigma]$ the function
$g=k\mu L(\sigma\lambda)/x_0<\frac12$ satisfies $g-g^2=k\mu Lh/x_0^2=c$. So $v=k\mu
L(\sigma\lambda)=v_0$ and $y=v/A=\mu y_0$.
\end{proof}

\subsection{\texorpdfstring{The coefficients $\hb_j$}{The second coefficient sequence}}
By \eqref{eq:balanced}, $\hb=t^{3/2}(t+G)^{-1}a$, where in the model $G=V_y^*V_y$ is the Hankel
matrix $G_{jk}=\int\lambda^{j+k}\abs y^2d\lambda$. The next lemma gives the lower bounds
required in (A4). For a positive Borel measure $\nu$ on $[0,\kappa]$ write
$G_\nu=(\int\lambda^{j+k}d\nu)_{j,k\ge0}$, a bounded positive operator on $\ell^2$.

\begin{lemma}\label{lem:two}
Let $0<w_-\le w_+$. There are $c_0,\eta_0>0$, depending only on $\kappa,w_-,w_+$, such that
the following holds. Let $\nu$ be a Borel measure on $[0,\kappa]$ with
$\nu\ge w_-\,d\lambda$ and $\nu([0,\kappa])\le w_+$, let $0<t\le1$, $0<\beta\le\frac12\kappa$,
and let $s\in\ell^2$ satisfy $\sup_j\abs{s_j-\beta^j/(j+1)}\le\eta_0$. Then
$e=(t+G_\nu)^{-1}s$ has $\abs{e_j}\ge c_0$ for at least two indices $j\ge2$.
\end{lemma}
\begin{proof}
Suppose not. Then there are $\nu_n,t_n,\beta_n,s^{(n)}$ as in the statement with
$\eta_0=1/n$, and $\theta_n\to0$, such that $e^{(n)}=(t_n+G_{\nu_n})^{-1}s^{(n)}$ has
$\abs{e^{(n)}_j}\le\theta_n$ for all $j\ge2$ except possibly one index $j_n\ge2$. Write $e$
for $e^{(n)}$ and $E(\lambda)=\sum_je_j\lambda^j$. The equation $(t_n+G_{\nu_n})e=s^{(n)}$
reads
\begin{equation}\label{eq:rows}
t_ne_j+\int\lambda^jE\,d\nu_n=s^{(n)}_j\qquad(j\ge0),
\end{equation}
and $E=e_0+e_1\lambda+e_{j_n}\lambda^{j_n}+R_n$ with $\sup\abs{R_n}\le\theta_n$. Let
$w_n=\int\lambda^{j_n}d\nu_n$ and let $\pi_n=\lambda^{j_n}\nu_n/w_n$, a probability measure.
Passing to a subsequence we may assume: $\nu_n\to\nu$ weak$^*$, where again
$\nu\ge w_-d\lambda$; $t_n\to t\in[0,1]$; $\beta_n\to\beta\in[0,\frac12\kappa]$, so that
$s^{(n)}_j\to\int\lambda^jd\pi^\beta$ for each $j$, where $\pi^\beta$ is the uniform
probability measure on $[0,\beta]$ ($\pi^0=\delta_0$); either $j_n=j_*$ for all $n$ or
$j_n\to\infty$; and $\pi_n\to\pi$ weak$^*$. In the first case $\pi=\lambda^{j_*}\nu/\!\int
\lambda^{j_*}d\nu$. In the second $\pi=\delta_\kappa$, because $\nu_n\ge w_-d\lambda$ near
$\kappa$ and $\nu_n$ has mass at most $w_+$.

Let $M_n=\abs{e_0}+\abs{e_1}+\abs{e_{j_n}}w_n$. Suppose first that $M_n$ is bounded; we may
assume $(e_0,e_1,e_{j_n}w_n)\to(\epsilon_0,\epsilon_1,\epsilon')$. Fix $j\ge2$ with
$j\ne j_*$. For large $n$, $\abs{e_j}\le\theta_n$, and letting $n\to\infty$ in
\eqref{eq:rows} gives $\int\lambda^jd\Xi=0$ for the signed measure
\[
\Xi=\pi^\beta-(\epsilon_0+\epsilon_1\lambda)\nu-\epsilon'\pi .
\]
So all moments of $\lambda^2\Xi$ vanish, except possibly the one of order $j_*-2$ when $j_*$
is finite. In either case $\lambda^{J}\Xi$ has all moments zero for some $J$, hence
$\lambda^J\Xi=0$ and $\Xi$ is carried by $\{0\}$. On $(\frac12\kappa,\kappa]$ the measure
$\pi^\beta$ vanishes. If $j_*$ is finite this gives $P\nu=0$ there for the polynomial
$P=\epsilon_0+\epsilon_1\lambda+\epsilon'\lambda^{j_*}/\!\int\lambda^{j_*}d\nu$, so $P=0$
because $\nu\ge w_-d\lambda$. If $j_n\to\infty$ it gives
$(\epsilon_0+\epsilon_1\lambda)\nu+\epsilon'\delta_\kappa=0$ on $(\frac12\kappa,\kappa]$, so
again $\epsilon_0=\epsilon_1=0$ and then $\epsilon'=0$. Thus
$\epsilon_0=\epsilon_1=\epsilon'=0$. But the limit of \eqref{eq:rows} for $j=0$ is
$t\epsilon_0+\int(\epsilon_0+\epsilon_1\lambda)d\nu+\epsilon'=1$, a contradiction.

If $M_n$ is unbounded, pass to a subsequence with $M_n\to\infty$, divide \eqref{eq:rows} by
$M_n$ and argue in the same way with
$\Xi=-(\epsilon_0+\epsilon_1\lambda)\nu-\epsilon'\pi$, where now
$\abs{\epsilon_0}+\abs{\epsilon_1}+\abs{\epsilon'}=1$: again $\Xi$ is carried by $\{0\}$,
which forces $\epsilon_0=\epsilon_1=\epsilon'=0$, a contradiction.
\end{proof}

\section{Failure of Claim (C)}

Theorem~\ref{thm:nogo} enters through the vectors $v_N^\alpha$.

\begin{lemma}\label{lem:osc}
Let $W=e^{i\alpha\sin(NT)}$, multiplication by $e^{i\alpha\sin N\lambda}$, and
$C=\frac12(W+W^*)$. Then $W$ and $W^*$ are power series in $T$ with absolute coefficient sums
at most $K_N=\exp(\alpha\sinh N)$, and
$\frac14(\norm{x_0}^2-\norm{Cx_0}^2)\to\frac18(1-J_0(2\alpha))$ as $N\to\infty$, uniformly in
$\alpha\in[0,\alpha_{\max}]$. For $\alpha\in[1.45,1.54]$ the limit is at least $0.15$.
\end{lemma}
\begin{proof}
$\sin(Nz)$ has power series coefficients of absolute sum $\sinh N$, and absolute sums of
coefficients are submultiplicative. $C$ is multiplication by $\cos(\alpha\sin N\lambda)$, so
$\norm{x_0}^2-\norm{Cx_0}^2=\int x_0^2\sin^2(\alpha\sin N\lambda)$, and the mean of
$\sin^2(\alpha\sin t)$ is $\frac12(1-J_0(2\alpha))$; use Fej\'er's lemma. Finally
$J_0\le-0.22$ on $[2.9,3.08]$.
\end{proof}

\begin{theorem}\label{thm:main}
Let $t_0$ be small, $x_0$ as in Lemma~\ref{lem:x0}, and $\alpha\in[1.45,1.54]$. There is
$\beta_*>0$, and for each $\beta\in(0,\beta_*]$ an $N_*(\beta)$, such that for
$N\ge N_*(\beta)$ the state $S(\beta,\alpha,N)=(y,a,t,v)$ has the following properties.
\begin{enumerate}[label=\textup{(\alph*)}]
\item $v$ is an admissible extremal image in the sense of
Definition~\ref{def:admissible}, with $0.46<\norm{x_0-v}<0.52$.
\item $\tau(v)\le2\mu^2t_0\,\beta/\sigma$. In particular $\tau(v)\to0$ as $\beta\to0$.
\item If $w\in H$ satisfies $\norm{w-y/\norm y}\le10^{-3}$ and $u$ is any extremal image with
\[
\ip{x_0-u}{w}=\ip{x_0-v}{w}\quad\text{and}\quad
\bigl|\norm{x_0-u}^2-\norm{x_0-v}^2\bigr|\le11\,\tau(v),
\]
then $\tau(u)\ge\frac1{10}\exp(-\alpha\sinh N)>0$.
\item More generally, every $u\in H$ with $\norm{u-v}\le\frac1{10}$ has
$\Delta(u)\ge\frac1{10}\exp(-\alpha\sinh N)$.
\end{enumerate}
In particular Claim~\textup{(C)} fails for $v'=v$, with $w=w_{00}$, although the defect of
$v'$ can be taken as small as we please. Part \textup{(d)} does not involve the bookkeeping
of \textup{(45)}: the ``room for variation'' that the constraint leaves, in the words of
\cite{E}, contains no vectors with defect below this fixed positive threshold (which
depends on the chosen state and tends to zero as $N$ increases).
\end{theorem}
\begin{proof}
Let $c_0,\eta_0$ be the constants of Lemma~\ref{lem:two} for $w_-=k^2/25$ and $w_+=2$; they
depend only on $\kappa$. Choose $\beta_*\le\sigma/400$, and $N_*(\beta)$ so large that
$\rhoN\le\min(10^{-5},\eta_0,\beta)$ and $\phi_N\le10^{-3}$ for $N\ge N_*(\beta)$, and the
quantity in Lemma~\ref{lem:osc} is at least $0.149$. All statements of
Proposition~\ref{prop:states} are then available.

(b) By Proposition~\ref{prop:states}(i), $t\le\tau(v)\le t(\Lambda(\beta)+\rhoN^2)\le2t$.

(a) (A1): $y$ is cyclic and lies in the cone by Proposition~\ref{prop:states}(iii). By (iv),
$\norm{x_0-v}^2$ is within $2\cdot10^{-3}$ of $\varepsilon_\infty^2(\beta,\alpha)$, which
differs from $\frac12(1-J_0(\alpha))\in[0.229,0.256]$ by at most $m(\beta)\le10^{-3}$.
(A2): by (ii), $\sum_{j\ge1}\abs{a_j}^2\le2\beta^2\le\frac1{100}\le\frac1{100}\abs{a_0}^2$.
(A3): by (i) and (ii), $\sum_{j\ge1}\abs{a_j}^2\ge\frac14\beta^2\ge\frac15\beta^2\norm a^2$.
Now $\beta=\sigma t/(\mu^2t_0)\ge\xi\tau(v)$ with
$\xi=\sigma/(2\mu^2t_0)=t_0/(2k^2\varepsilon_0^2)$, and $\tau(v)\le2t_0$, so
$\frac15\beta^2\ge\frac15\xi^2\tau^2\ge\tau^{15}$ once $t_0$ is small.
(A4): for $j\in\{2,3,4,5\}$, $\abs{\ha_j}=t^{1/2}\abs{a_j}\ge t^{1/2}\beta^5/6\ge
(\tau/2)^{1/2}(\xi\tau)^5/6>\tau^{30}$ once $t_0$ is small. For $\hb$ apply
Lemma~\ref{lem:two} with $\nu=\abs y^2d\lambda$ and $s=a$: its hypotheses hold by
Proposition~\ref{prop:states}(i),(iii), so $\abs{\hb_j}=t^{3/2}\abs{e_j}\ge c_0t^{3/2}\ge
c_0(\tau/2)^{3/2}>\tau^{30}$ for two indices $j\ge2$, once $t_0$ is small. Choose $j_3,j_4$
to be these and $j_1,j_2\in\{2,3,4,5\}$ different from them.

(d) By
Proposition~\ref{prop:states}(v), $\norm{v-v_N^\alpha}^2\le10^{-3}$, so
$\norm{u-v_N^\alpha}^2\le(0.1+0.032)^2\le0.02$, and Theorem~\ref{thm:nogo}(ii) with
Lemma~\ref{lem:osc} gives $\Delta(u)\ge K_N^{-1}(0.149-0.02)$.

(c) Put $y'=y/\norm y$. By Proposition~\ref{prop:states}(iii),
$\norm{v/\norm v-y'}\le10^{-3}$, so Lemma~\ref{lem:confine} with $\vartheta=10^{-3}$ gives
$\norm{u-v}^2\le13\tau(v)+0.008\le0.01$, using (b) and $\beta\le\sigma/400$. Now apply (d)
and $\Delta(u)=\tau(u)$.
\end{proof}

\section{The states are reachable}\label{sec:reach}

The argument of \cite{E} for Claim (C) makes no use of how the state $v'$ was reached, so
Theorem~\ref{thm:main} already shows that it is not valid. One may still ask whether the
bookkeeping of the MC could keep the construction away from such states. The MC proceeds by
steps
\begin{equation}\label{eq:step}
y_{n+1}=y_n+\sum_{j\ge0}r_jT^jy_n
\end{equation}
\textup{(\cite{E}, (13))}, chosen so that the defect decreases \textup{(31)}, the quantities
$\norm{x_0-v}^2$ and $\ip{Rx_0}{x_0}$ change by at most a multiple of the decrease
\textup{(32), (45)}, and the leading coefficient stays dominant with the tail bounded below
\textup{(33)}. Apart from the remark that the degrees $j$ involved are bounded by a function
of the defect \textup{(after (34))}, no further rule is given. The next theorem shows that the
requirements \textup{(31)--(33)} and \textup{(45)} do not exclude the states of
Theorem~\ref{thm:main}.

\begin{theorem}\label{thm:path}
Let $t_0$ be small and $x_0$ as in Lemma~\ref{lem:x0}. For every $\eta>0$ and $\tau_*>0$ there
are $n$ and vectors $y^0,y^1,\dots,y^n$, with extremal images $v^i$ of $y^i$ and defects
$\tau_i=\tau(v^i)$, such that:
\begin{enumerate}[label=\textup{(\alph*)}]
\item $y^0=y_0=\frac{\sqrt3}2u_0$ and $v^0=v_0$, the extremal image of Lemma~\ref{lem:x0};
\item each step is of the form \eqref{eq:step}:
$y^{i+1}=y^i+\sum_jr^{(i)}_jT^jy^i$ with $\sum_j\abs{r^{(i)}_j}<\infty$; and
$\norm{y^{i+1}-y^i}\le\eta$, $\norm{v^{i+1}-v^i}\le\eta$;
\item every $v^i$ is admissible;
\item $\tau_0>\tau_1>\dots>\tau_n$ and $\tau_n<\tau_*$;
\item $\bigl|\norm{x_0-v^i}^2-\norm{x_0-v^0}^2\bigr|\le\eta\tau_0$ for all $i$, and
$\bigl|\norm{x_0-v^{i+1}}^2-\norm{x_0-v^i}^2\bigr|\le\tau_i-\tau_{i+1}$ for each step; so
in each step $\norm{x_0-v}^2$ and $\ip{Rx_0}{x_0}=\norm{x_0-v}^2+\tau(v)$ change by at most
twice the decrease of the defect, and $\ip{Rx_0}{x_0}$ varies by at most $2\tau_0$ along the
chain;
\item the conclusion of Theorem~\ref{thm:main}(c) holds for $(y,v)=(y^n,v^n)$, with
$\frac1{20}$ in place of $\frac1{10}$ and some $\alpha\in[1.45,1.54]$, $N\in\N$.
\end{enumerate}
\end{theorem}
\begin{proof}
\emph{Step 1: a path of explicit states.} We first check that $D(\beta)\le D(\sigma)
=\varepsilon_0^2$ for $\beta\le\sigma$. On $[0,\sigma]$ we have $x_0(1-g_\sigma)=h$ by
Proposition~\ref{prop:states}. As $L$ is increasing, $0\le g_\sigma-g_\beta\le
6t_0\,\sigma(\sigma-\beta)$ on $[0,\beta]$, so $x_0(1-g_\beta)\le h(1+7t_0\sigma(\sigma-\beta))$
and $D(\beta)\le\beta h^2(1+15t_0\sigma(\sigma-\beta))\le\sigma h^2$. Also
$D(\beta)\ge(1-12t_0)m(\beta)$ and $m(\beta)\le\frac12$. Hence
\[
\Phi(\beta):=\frac{\varepsilon_0^2-D(\beta)}{1-m(\beta)}\in\bigl[0,\ \varepsilon_0^2+12t_0\bigr]
\subset[0,0.253],
\]
$\Phi$ is continuous, $\Phi(\sigma)=0$ and $\Phi(\beta)\to\varepsilon_0^2$ as $\beta\to0$.
Since $J_0$ is strictly decreasing on $[0,\alpha_{\max}]$ with
$\frac12(1-J_0(\alpha_{\max}))>0.255$, there is a unique continuous
$\alpha:(0,\sigma]\to[0,\alpha_{\max})$ with
$\frac12(1-J_0(\alpha(\beta)))=\Phi(\beta)$. Then $\alpha(\sigma)=0$,
\[
\varepsilon_\infty^2(\beta,\alpha(\beta))=\varepsilon_0^2\quad\text{for all }\beta ,
\]
and $\alpha(\beta)\to\alpha_1$ as $\beta\to0$, where $J_0(\alpha_1)=1-2\varepsilon_0^2$; as
$\varepsilon_0^2=\frac14+O(t_0^2)$, $\alpha_1\in(1.5,1.54)$.

Fix $\beta_1\le\beta_*$ so small that $2t(\beta_1)<\tau_*$ and
$\alpha(\beta_1)\in[1.45,1.54]$, and consider the states
\[
S_\beta=S(\beta,\alpha(\beta),N),\qquad\beta_1\le\beta\le\sigma ,
\]
for a large $N$ to be chosen. Since $t(\beta)\ge t(\beta_1)$ we have
$\rhoN\le\omega_N/t(\beta_1)$ on the whole path. By Proposition~\ref{prop:states}: $S_\sigma$
is the state $(\mu y_0,v_0)$; the path is norm-continuous; its defects are
$\tau(\beta)=t(\beta)(\Lambda(\beta)+\delta)$ with $0\le\delta\le\rhoN^2$, where
$t(\beta)\Lambda(\beta)$ is strictly increasing in $\beta$; and
$\bigl|\norm{x_0-v}^2-\varepsilon_0^2\bigr|\le100\rhoN+\phi_N$ along it. Each $S_\beta$ is
admissible for $N$ large: (A1) by Proposition~\ref{prop:states}(iii) and because
$\norm{x_0-v}^2$ is close to $\varepsilon_0^2$; (A2)--(A4) by the computations in the proof of
Theorem~\ref{thm:main}, which used only $\beta\le\sigma$, $\tau\le2t_0$, and
$\rhoN\le\min(\eta_0,\beta)$.

By the comparisons in the proof of Proposition~\ref{prop:states}, up to an error $O(\rhoN)$
in norm the states are given by $g_\beta$ on $[0,\beta]$ and by $g_N$ on $(\beta,\kappa]$, and
$\abs{e^{i\alpha s}-e^{i\alpha's}}\le\abs{\alpha-\alpha'}$. So there is a mesh
$\sigma=\beta^0>\beta^1>\dots>\beta^n=\beta_1$, independent of $N$, such that for all large
$N$ consecutive states differ by at most $\eta/4$ in $y$ and in $v$. Let $\gamma>0$ be the
smallest of the gaps $t(\beta^i)\Lambda(\beta^i)-t(\beta^{i+1})\Lambda(\beta^{i+1})$. Then
choose $N\ge N_*(\beta_1)$ so large that in addition $t(\beta)\rhoN^2<\gamma/4$ and
$100\rhoN+\phi_N<\min(\gamma/8,\eta\tau_0/2)$ on the path; then consecutive defects differ by
more than $\gamma/2$ and consecutive values of $\norm{x_0-v}^2$ by less than $\gamma/4$. The states $S_{\beta^i}$ then have properties (c)--(f), the last
with constant $\frac1{10}$, and the first of them is $(\mu y_0,v_0)$.

\emph{Step 2: steps of the form \eqref{eq:step}.} Write $S_{\beta^i}=(\mu\tilde y^i,\cdot,
\cdot,\tilde v^i)$, so $\tilde y^0=y_0$. By Proposition~\ref{prop:states} each
$F_i=\tilde y^i/y_0$ is a bounded measurable function, bounded away from $0$, with
$\re F_i>0$. So $\log F_i$ is bounded, and there are polynomials $q_{i,m}$, bounded on
$[0,\kappa]$ uniformly in $m$, with $q_{i,m}\to\log F_i$ a.e.\ as $m\to\infty$; hence
$e^{q_{i,m}(T)}y_0\to\tilde y^i$ in norm. Take $q_{0,m}=0$. The vector
$\tilde y^i$ has the same extremal image as $\mu\tilde y^i$, with Tikhonov parameter
$t(\beta^i)/\mu^2$. Now all the properties in (c)--(f)
are stable under small perturbations of the vectors $\tilde y^i$ with these parameters held
fixed: for fixed $t>0$ the Tikhonov data $a=(V_y^*V_y+t)^{-1}V_y^*x_0$, $v=V_ya$ and $\hb$
depend norm-continuously on $y$, the conditions (A1)--(A4) and (d), (e) are finitely many
strict inequalities between continuous functions of these data, and (f) was derived from
strict inequalities for $\norm{v-v_N^\alpha}$, $\norm{v/\norm v-y/\norm y}$ and $\tau(v)$
with room to spare. So for $m$ large the vectors $y^i=e^{q_{i,m}(T)}y_0$ have properties
(a) and (c)--(f), each $y^i$ is non-zero a.e.\ and hence cyclic, and consecutive vectors
are within $\eta$. Finally
\[
y^{i+1}=e^{(q_{i+1,m}-q_{i,m})(T)}y^i ,
\]
and $e^{q(T)}=I+\sum_{j\ge0}r_jT^j$ with $\sum_j\abs{r_j}<\infty$ for every polynomial $q$.
\end{proof}

\begin{remark}\label{rem:path}
(1) The chain starts at the vector $y_0$ of \cite{E} and at an admissible extremal image of
it. Along it the defect decreases, the states are admissible, and the distance to $x_0$ is
constant to within any tolerance, both in total and relative to the decrease of the defect in
each step; these are the requirements \textup{(31)--(33)} and \textup{(45)} of \cite{E}, the
upper bound in \textup{(33)} being understood as in the remarks after
Definition~\ref{def:admissible}. The chain ends at a state where Claim (C) fails.

(2) The vectors $y^i$ are of the form $F(T)y_0$ with $F$ entire and zero-free on
$[0,\kappa]$. They approximate functions with a moving jump, and the coefficients
$r^{(i)}_j$ are large even though $\norm{y^{i+1}-y^i}$ is small. In particular the
degrees of the steps are not bounded in terms of the defect, as \cite{E} envisages after its
(34). The formulas of \cite{E} for the effect of a step are first-order expansions without
remainder estimates, and we do not claim that they are accurate for these steps. What
Theorem~\ref{thm:path} shows is that the conditions which \cite{E} imposes on the
\emph{states} of the MC do not distinguish this chain from a successful run; whether a
specific rule for choosing small steps would avoid such states is not something the text of
\cite{E} allows one to decide.

(3) In the model, $v_0$ lies within $O(t_0)$ of the genuine solution
$x_0\one_{(\sigma,\kappa]}$, which has defect $0$. The chain moves away from it, to a state
at distance about $\frac12$ from every vector of defect $0$, while its defect decreases
monotonically.
\end{remark}

\section{Discussion}\label{sec:discussion}

\subsection*{What this shows}
The convergence step of \cite{E} rests on an implication of the form: \emph{type 1 and an
admissible state with small defect together guarantee that constrained descent steps exist
all the way to defect $0$.} The model operator satisfies the hypotheses on the operator, the
states of Theorem~\ref{thm:main} satisfy the hypotheses on the state, and the conclusion is
false. So the implication is not valid, and with it the proof of norm convergence. Since one
branch of the type 2 case of \cite{E} ends by invoking the type 1 construction, the argument
as a whole is incomplete at its central point.

The mechanism is simple. Small defect does not localise: multiplying part of $x_0$ by a
rapidly oscillating unimodular function makes every pairing $\ip{T^jv}{x_0-v}$ small by
cancellation, while $v$ stays far from every vector of defect $0$ (which in the model are the
vectors $x_0\one_S$). The quantities that \cite{E} controls are all of this averaged kind
and cannot detect the oscillation. And by Lemma~\ref{lem:confine}, the confinement constraint
can only succeed if the small ball it defines already contains vectors of arbitrarily small
defect.

\subsection*{Why the test is made in an operator with invariant subspaces}
Theorem~\ref{thm:nogo}(ii) needs a non-scalar unitary $W$ such that $W$ and $W^*$ are both
absolutely convergent power series in $T$. A self-adjoint operator has many. An operator
without invariant subspaces has none: such a $W$ commutes with $T$, and the spectral
subspaces of a non-scalar unitary commuting with $T$ are invariant under $T$. So our
states do not exist for a hypothetical counterexample to the invariant subspace problem, and
we have no information on what the construction would do there.

This limitation cannot be removed by a better example. In \cite{E}, Claim (C) is applied
repeatedly, with constraints $w_{00},w_{01},\dots$ that confine the states to smaller and
smaller balls (Lemma~\ref{lem:confine}), and the outcome is a norm-convergent sequence with
defect tending to $0$, hence a non-cyclic vector. For an operator without invariant subspaces
this iterated form of the claim is therefore false, and it is the statement from which the
contradiction is drawn; there is no independent way to examine it in such an operator. A proof
by contradiction can only be checked by checking the derivation, and a derivation is tested by
running it in a setting where everything it uses is true. The type 1 argument of \cite{E} uses the absence of
invariant subspaces in the following places only, and each holds in the model for the vectors
concerned:
\begin{itemize}[leftmargin=2em]
\item the standing reductions, that $T$ is injective with dense range and is not
invertible (here: Lemma~\ref{lem:model}(i));
\item the vectors $y_n$ to which the extremal construction is applied are cyclic (here: by
Proposition~\ref{prop:states} and Theorem~\ref{thm:path});
\item in the proof of its Lemma~2, an equation $dy=\sum_{j\ge1}c_jT^jy$ with $d\ne0$,
$\abs{c_j}\le1$ is impossible for every $y\ne0$, in particular for weak limits of the vectors
considered (here: Lemma~\ref{lem:model}(iii), which holds for all $y\ne0$);
\item the type 1 inequalities (here: Lemma~\ref{lem:model}(ii)).
\end{itemize}
The remaining appeals are of the form ``or else we already have an invariant subspace'',
which are exits, not hypotheses. So if the derivation of Claim (C) were sound, it would apply
verbatim in the model, where its conclusion fails. What our result leaves open is whether
some \emph{other} argument, using the absence of invariant subspaces in an essential and
quantitative way at this step, could establish continuation. The text of \cite{E} contains no
such argument, and since in that setting continuation is what yields the theorem, it would
have to carry the whole weight of the proof.

\subsection*{Other problems}
Independently of the convergence step, the Main Construction of \cite{E} is specified only to
first order, with no control of step sizes or of higher-order terms, and the estimate (27) of
\cite{E} used to prepare its starting state is incorrect. These are discussed in detailed notes
accompanying this paper, which also give corrected proofs of the parts of \cite{E} that are
sound.

\subsection*{Formalisation and numerical checks}
Theorem~\ref{thm:nogo} and the identity in Lemma~\ref{lem:confine} are machine-checked in
Lean~4 with Mathlib. The results of Sections~\ref{sec:model}--\ref{sec:reach} are proved here
by hand and are not formalised. The script \texttt{experiments/explicit\_states.py} checks on
a discretisation that the series \eqref{eq:aj} agrees with the numerically computed extremal
series (to about $10^{-14}$); this tests the identity of Lemma~\ref{lem:explicit}, at
frequencies $N$ far below those the estimates require. It also tabulates the defect, distance,
tail and the vector $e=t^{-3/2}\hb$ along the path of Theorem~\ref{thm:path} in the limit
$N\to\infty$.

\begin{thebibliography}{9}
\bibitem{E} P.~H.~Enflo, \emph{On the invariant subspace problem in Hilbert spaces},
arXiv:2305.15442v2 (2024).
\end{thebibliography}

\end{document}
